% Chapter 14 -- Working with Several Applications (Part II, Reference).
% Sources (Law 1): os/milton/kmain.c (the co-resident tenants: FINDER,
% NOTES and HELLO launched at start-up, HELLO foreground;
% flair_live_sync_foreground: the front window's owner is the foreground
% application and its bar is band 2; a click on a Finder window or the
% desktop makes the Finder foreground; finder_power / finder_exit_dos: the
% open disk application is asked to quit first, FINDER-POWER-CANCEL when it
% refuses; the cooperative pump: each tenant runs until it returns);
% os/flair/tbxgate.c (ONE disk-application slot: a second launch is refused
% with TENANT-SLOT-BUSY; the Calculator and every .EXE share it; a tenant
% that faults is closed, FLAIR-TENANT-EXIT); os/flair/process.c (the
% process table); os/flair/scrap.c (one desk Scrap for every tenant).
% ADR-0013 (the App Contract: co-resident tenants, one disk slot).
% Captures: screens/m_*.mouse.  Model: System 6 User's Guide (1988)
% Ch. 10 "MultiFinder" and the Application menu topic of What's New in
% System 7 (apple-macintosh.md Sec 8.1, 8.3), cut to what FLAIR ships.
% GATES: test-flair-appswitch test-flair-tenant-menu test-flair-app-launch test-tenant-load test-flair-close-terminate test-flair-shutdown test-flair-exit-to-dos test-scrap test-pubs-flair-screens
\Gates{test-flair-appswitch, test-flair-tenant-menu, test-flair-app-launch,
  test-tenant-load, test-flair-close-terminate, test-flair-shutdown,
  test-flair-exit-to-dos, test-scrap, test-pubs-flair-screens}
\chapter{Working with Several Applications}
\label{ch:multi}
\index{applications!several open}

FLAIR keeps more than one application open at a time.  This chapter
describes which applications can be open together, how you move from one to
another, and how they share the computer.

\section{What can be open at once}
\begin{center}\small
\begin{tabular}{@{}p{32mm}p{67mm}@{}}
\textbf{Application} & \textbf{When it is open}\\\hline
The Finder & always, from the moment FLAIR starts\\
HELLO & from start-up until you close its window\\
NOTES & from start-up until you close its window\\
One application from the disk & from the moment you open it until it
quits: InitechText, InitechCalc, InitechPaint, the Calculator, Initech
123, InitechWord or a sample application\\
\end{tabular}
\end{center}
All of them can be open together, each in its own windows.  Only one is
the \term{active application}: the one whose window is in front, and whose
menus are in the lower menu bar.

\section{Moving from one application to another}
\index{applications!switching between}\index{active application}
To make an application active, click one of its windows.  To make the
Finder active, click a Finder window or a bare part of the desktop.

\begin{TaskSteps}
\Step With InitechText open, click a part of a Finder window that shows
beside InitechText's window.
\Result{The Finder becomes the active application.  Its windows come to
the front and its menus fill the lower menu bar.  InitechText's window
stays open behind them; its title is grey to show that it is not active.}
\Step Choose \menuitem{Close Window} from the \menuname{File} menu, or press
\key{Ctrl-W}, to close the front Finder window.
\Result{The next Finder window is now in front.}
\end{TaskSteps}

\FlairScreen[label=fig:mswitch]{m_switch}{The Finder active, InitechText's
window behind it}

\begin{TaskSteps}[resume]
\Step Click InitechText's window.
\Result{InitechText becomes the active application again, with its
document as you left it.  The first click only brings the window to the
front; click again to place the insertion point.}
\end{TaskSteps}

An application that is not active keeps its windows and their contents.
Nothing is lost by switching, and you do not need to save before you click
another window --- although it is good practice.

When an application's window is hidden behind others, move or collapse the
windows in front of it, or close the Finder windows you do not need, until
you can click it.

\section{One application from the disk at a time}
\index{applications!one at a time}
\Operation{TENANT-SLOT-BUSY}
FLAIR has room for one application from the disk.  While one is open,
opening another does nothing: double-clicking its icon, or choosing the
Calculator from the Apple menu, leaves the screen as it was.

\FlairScreen[label=fig:mcalcbusy]{m_calcbusy}{The Calculator chosen from the
Apple menu while InitechText is open: it does not open}

To use the other application, quit the open one first: choose
\menuitem{Quit} from its \menuname{File} menu, or close its window.  HELLO
and NOTES are part of FLAIR and do not count: they stay open beside the
application from the disk.

\section{Sharing the computer}
\index{cooperative operation}
The open applications take turns.  When you click, type or choose a
command, FLAIR passes it to the active application, which does the work and
hands the computer back.  While an application is busy --- InitechCalc
while its \texttt{FILE COPY} box is on the screen, for example --- the
desktop waits: clicks and keys are ignored until the work is done.

If an application stops working, FLAIR closes it: its windows disappear and
its menus leave the menu bar.  The Finder, HELLO and NOTES go on working.
Work the application had not saved is lost.

\section{Copying between applications}
\index{Clipboard!between applications}
There is one Clipboard for all applications.  Text you cut or copy in one
application can be pasted in another, or in a later document of the same
application, until you cut or copy something else or restart the computer.
In this version of FLAIR, InitechText is the application that uses the
Clipboard (Chapter~\ref{ch:appsref}).

\section{Quitting everything at once}
\index{quitting!all applications}
\Operation{FINDER-POWER}
\menuitem{Restart}, \menuitem{Shut Down} and \menuitem{Exit to InitechDOS}
first ask the open application from the disk to quit, just as if you had
chosen its \menuitem{Quit} command.  An application that holds unsaved
work asks you what to do with it.

\FlairScreen[label=fig:mshutask]{m_shutask}{Shut Down, with unsaved work in
InitechText}

\SideHead{Save}The application saves the work and quits; then the command
goes on.
\SideHead{Don't Save}The application quits without saving; then the
command goes on.
\SideHead{Cancel}The application stays open with its work, and the command
stops: the computer is not shut down or restarted, and FLAIR is not left.

HELLO and NOTES hold nothing that needs saving and are closed with FLAIR.
